\section{Limitations, reproducibility and resources}
\label{sec:limits}

\paragraph{Main limitations.}
\begin{itemize}[leftmargin=*]
  \item \emph{Scope.} One model (Qwen2.5-3B-Instruct), synthetic city questions in three task templates, one
        instruction and note format, and one agent run. The planned reworded-instruction split was dropped before the
        first freeze, so transfer to new instructions is untested.
  \item \emph{Unresolved comparison.} The comparison with logit bias is inconclusive (recovery difference 95\% CI
        $-0.043$ to $+0.088$); neither superiority nor equivalence was established at the frozen margins.
  \item \emph{Costs of the repair.} The edit broke previously correct answers (base-weighted damage 0.030, upper bound
        0.055), while B1 and B4 broke none. Counterfactual following (0.917) cleared its 0.90 threshold by a small
        margin. Both thresholds were applied to point estimates.
  \item \emph{Dose.} The selected dose was the largest in the frozen grid, and larger doses were not searched.
  \item \emph{Controls.} The only random-edit control in confirmation was B6, five size-matched perturbations of the
        same weight entries, used in the pathway tests. Random-head and random-direction controls belong to the
        follow-up. The prompt offers ``unknown'', and the present results do not distinguish a change in how the
        model treats that option from a change in how it uses the evidence.
  \item \emph{Mechanism tests.} The account test faced a predictor ceiling (baseline AUROC 0.998) and fitted the
        agent's account poorly (\cref{sec:explain}). The pathway result is sufficiency, not necessity.
  \item \emph{The note.} The note is a prompted summary. We do not treat it as a faithful trace of the model's
        reasoning.
\end{itemize}

\paragraph{Reproducibility.}
The design was frozen in three stages before the data each stage governed were generated: F0 (data, grader,
comparators, edit rules), F1 (bounds, run procedure, cost limits) and F2 (candidates, doses, analysis code). All
model outputs on H and T-city were generated after F2. Besides two deviations recorded at F1, four tooling fixes
were made: a comparator file missing from the run bundles, smoke tests that could not apply full-size edits, a
missing key that crashed the first selection run (\$0.09, counted in the cost below), and a guard failure that
exited without a log message. None changed a scoring, selection or decision rule, and each was reviewed and accepted
in the project record. \Cref{app:artefacts} lists the commits, hashes, code and per-row outputs, and states which
artefacts are not yet public.

\paragraph{Compute and cost.}
The agent run cost \$3.58 in API usage. Replaying the submissions on E took about 59 hours on a laptop CPU at no
cloud cost. Selection on S ran on one NVIDIA A40 GPU and cost \$4.38, including \$0.09 for a failed first attempt;
confirmation on H and T-city took 4.36 hours on one A40 and cost \$2.57. Total GPU charges were \$6.95.
